\sectiontoc{Introduction}

Perturbations of a spatially isotropic and homogeneous expanding 
universe have been investigated in a Newtonian approximation by 
Bonnor~\cite{Bonnor1957} and relativistically by Lifshitz~\cite{Lifshitz1946}, 
Lifshitz and Khalatnikov~\cite{LifshitzKhalatnikov1963}, 
and Irvine~\cite{Irvine1965}. Their method was to consider small 
variations of the metric tensor. This has the disadvantage that the 
metric tensor is not a physically significant quantity, since one 
cannot directly measure it, but only its second derivatives. 
It is thus not obvious what the physical interpretation of a given 
perturbation of the metric is. Indeed it need have no physical 
significance at all, but merely correspond to a coordinate transformation. 
Instead it seems preferable to deal in terms of perturbations of the 
physically significant quantity, the curvature.
\sectiontoc{Notation}

Space--time is represented as a four--dimensional Riemannian space 
with metric tensor \(g_{ab}\) of signature \(+2\). Covariant differentiation 
in this space is indicated by a semi--colon. Square brackets around 
indices indicate antisymmetrisation and round brackets symmetrisation.  
The conventions for the Riemann and Ricci tensors are:
\[
\nabla_{a;[bc]} = 2\, R^{p}{}_{a c b}\, \nabla_{p},
\]
\[
R_{ab} = R^{p}{}_{a p b}.
\]

\(\eta_{abcd}\) is the alternating tensor.  

Units are such that \(\kappa\), the gravitational constant, and \(c\), the speed of light, are one


\sectiontoc{The Field Equations}

We assume the Einstein equations:
\[
R_{ab} - \tfrac{1}{2} g_{ab} R = -T_{ab}
\]
where \(T_{ab}\) is the energy--momentum tensor of matter.  
We will assume that the matter consists of a perfect fluid. Then,
\[
T_{ab} = \mu\, u_{a} u_{b} + p\, h_{ab}
\]
where \(u_{a}\) is the velocity of the fluid, \(u_{a} u^{a} = -1\) :
\[
\mu \ \text{is the density}, \qquad 
p \ \text{is the pressure},
\]
\[
h_{ab} = g_{ab} + u_{a} u_{b} \quad 
\text{is the projection operator into the hyperplane orthogonal to } u_{a},
\]
\[
h_{ab} u^{b} = 0.
\]

\medskip

We decompose the gradient of the velocity vector \(u_{a}\) as
\[
u_{a;b} = \omega_{ab} + \sigma_{ab} + \tfrac{1}{3} h_{ab} \, \theta - \dot{u}_{a} u_{b}
\]
where
\[
\dot{u}_{a} = u_{a;b} u^{b} \quad \text{is the acceleration},
\]
\[
\theta = u^{a}{}_{;a} \quad \text{is the expansion},
\]
\[
\sigma_{ab} = u_{(c;d)} h^{c}{}_{a} h^{d}{}_{b} - \tfrac{1}{3} h_{ab} \theta 
\quad \text{is the shear},
\]
\[
\omega_{ab} = u_{[c;d]} h^{c}{}_{a} h^{d}{}_{b} 
\quad \text{is the rotation of the flow lines } u^{a}.
\]

We define the rotation vector \(\omega_{a}\) as
\[
\omega_{a} = \tfrac{1}{2} \eta_{abcd} \, \omega^{cd} u^{b}.
\]

\medskip

We may decompose the Riemann tensor \(R_{abcd}\) into the Ricci tensor 
\(R_{ab}\) and the Weyl tensor \(C_{abcd}\):
\[
R_{abcd} = C_{abcd} - g_{a[c} R_{d]b} + g_{b[c} R_{d]a} - \tfrac{1}{3} R g_{a[c} g_{d]b},
\]
\[
C_{abcd} = C_{[ab][cd]}, \qquad C^{a}{}_{bca} = 0 = C_{a[ bcd ]}.
\]

\(C_{abcd}\) is that part of the curvature that is not determined locally by 
the matter. It may thus be taken as representing the free gravitational 
field (Jordan, Ehlers and Kundt\cite{JordanKundt1960}).  
We may decompose it into its ``electric'' and ``magnetic'' components:
\[
E_{ab} = - C_{a b p q}\, u^{p} u^{q}, 
\]
\[
H_{ab} = - \tfrac{1}{2} C_{a p q r}\, \eta^{qr}{}_{bs}\, u^{p} u^{s},
\]
\[
C_{ab}{}^{cd} 
= 8 u_{[a} E_{b]}{}^{[c} u^{d]} 
+ 4 \delta_{[a}{}^{[c} E_{b]}{}^{d]}
- 2 \eta_{ab}{}^{cd} u^{p} H_{p}{}^{[c} u^{d]}
- 2 \eta^{cdrs} u_{r} H_{s[a} u_{b]} .
\]

\[
E_{ab} = E_{(ab)}, 
\qquad H_{ab} = H_{(ab)},
\]
\[
E_{a}{}^{a} = H_{a}{}^{a} = 0,
\]
\[
E_{ab} u^{b} = H_{ab} u^{b} = 0.
\]

\(E_{ab}\) and \(H_{ab}\) each have five independent components.  

We regard the Bianchi identities,
\[
R_{ab[cd;e]} = 0,
\]
as field equations for the free gravitational field.  

Then
\[
C_{abc d}{}^{d} = - R_{c[b;a]} + \tfrac{1}{6} g_{c[b} R_{;a]},
\]
(Kundt and Trümper\cite{KundtTrumper1962}).

Using the decompositions given above, we may write these in a form 
analogous to the Maxwell equations:
\begin{align}
h^{b}{}_{e} H_{bc;d} h^{cf} + 3 H_{ab} \omega^{b} - \eta_{abcd} u^{b} \sigma^{c}{}_{e} H^{de} 
&= \tfrac{3}{2} h_{a}{}^{b} \mu_{;b}, \tag{1} \\[1ex]
h^{b}{}_{e} E_{bc;d} h^{cf} + 3 E_{ab} \omega^{b} - \eta_{abcd} u^{b} \sigma^{c}{}_{e} E^{de} 
&= (\mu + p) \, \omega_{a}, \tag{2} \\[1ex]
\perp \dot{E}_{ab} + h^{c}{}_{(a} \eta_{b)cde} u^{c} H^{de} + E_{ab} \theta - E^{c}{}_{(a} \omega_{b)c} 
&- E^{c}{}_{(a} \sigma_{b)c} - \eta_{cde(a} \eta_{bpqr} u^{c} u^{p} \sigma^{d}{}_{q} E^{er} \notag \\ 
&+ 2 H^{c}{}_{(a} \eta_{b)cde} u^{c} \dot{u}^{e} = -\tfrac{1}{2} (\mu + p) \sigma_{ab}, \tag{3} \\[1ex]
\perp \dot{H}_{ab} - h^{c}{}_{(a} \eta_{b)cde} u^{c} E^{de} + H_{ab} \theta - H^{c}{}_{(a} \omega_{b)c} 
&- H^{c}{}_{(a} \sigma_{b)c} - \eta_{cde(a} \eta_{bpqr} u^{c} u^{p} \sigma^{d}{}_{q} H^{er} \notag \\ 
&+ 2 H^{c}{}_{(a} \eta_{b)cde} u^{c} \dot{u}^{e} = 0. \tag{4}
\end{align}

where \(\perp\) indicates projection by \(h_{ab}\) orthogonal to \(u_{a}\) (c.f.\ Trümper (7)).  

\medskip

The contracted Bianchi identities give,
\begin{align}
(R_{ab} - \tfrac{1}{2} g_{ab} R)^{;b} &= - T_{ab}{}^{;b} = 0, \notag \\
\dot{\mu} + (\mu + p) \theta &= 0, \tag{5} \\
(\mu + p) \dot{u}_{a} + p_{,b} h^{b}{}_{a} &= 0. \tag{6}
\end{align}

The definition of the Riemann tensor is,
\[
u_{a;[bc]} = 2 R_{a p b c} u^{p}.
\]

Using the decompositions as above we may obtain what may be regarded 
as ``equations of motion''.

\begin{align}
\dot{\theta} &= 2 \omega^{2} - 2 \sigma^{2} - \tfrac{1}{3} \theta^{2} + \dot{u}^{a}{}_{;a} 
- \tfrac{1}{2} (\mu + 3p), \tag{7} \\[1ex]
\perp \dot{\omega}_{ab} &= - \tfrac{2}{3} \omega_{ab} \theta 
+ 2 \sigma_{c[a} \omega_{b]}{}^{c} 
+ \dot{u}_{[p;q]} h^{p}{}_{a} h^{q}{}_{b}, \tag{8} \\[1ex]
\perp \dot{\sigma}_{ab} &= E_{ab} - \omega_{ac} \omega^{c}{}_{b} - \sigma_{ac} \sigma^{c}{}_{b} 
- \tfrac{2}{3} \sigma_{ab} \theta \notag \\
&\quad - \tfrac{1}{3} h_{ab} (2 \omega^{2} - 2 \sigma^{2} + \dot{u}^{c}{}_{;c}) 
+ \dot{u}_{a} \dot{u}_{b} 
+ \dot{u}_{(p;q)} h^{p}{}_{a} h^{q}{}_{b}, \tag{9}
\end{align}

where \(\; 2\omega^{2} = \omega_{ab} \omega^{ab}, \quad 2\sigma^{2} = \sigma_{ab} \sigma^{ab}\).

\medskip

We also obtain what may be regarded as equations of constraint:
\begin{align}
\theta_{,b} h^{ba} &= \tfrac{3}{2} \Big[ 
  (\omega_{bc;b} + \sigma^{c}{}_{b;c}) h^{ca} 
  - \dot{u}^{b} (\omega_{ab} + \sigma_{ab}) \Big], \tag{10} \\[1ex]
\omega^{a}{}_{;a} &= 2 \omega_{a} \dot{u}^{a}, \tag{11} \\[1ex]
H_{ab} &= - h^{c}{}_{(a} \eta_{b)cde} u^{c} 
\left( \omega^{d;e} + \sigma^{d;e} \right). \tag{12}
\end{align}

\medskip

We consider perturbations of a universe that in the undisturbed state 
is conformally flat, that is
\[
C_{abcd} = 0.
\]

By equations (1)–(3), this implies
\[
\sigma_{ab} = \omega_{ab} = 0, 
\qquad h^{b}{}_{a} \mu_{;b} = 0, 
\qquad \theta_{;b} h^{ba}.
\]
If we assume an equation of state of the form, $\mu = \mu(\rho)$,  
then by (6), (10),
\[
\mu_{,b} h^{b}{}_{a} = 0 = \dot{\mu} u_{a}.
\]

This implies that the universe is spatially homogeneous and isotropic 
since there is no direction defined in the 3-space orthogonal to $u_{a}$.

In this universe we consider small perturbations of the motion 
of the fluid and of the Weyl tensor. We neglect products of small 
quantities and perform derivatives with respect to the undisturbed 
metric. Since all the quantities we are interested in with the 
exception of the scalars, $\mu, \dot{\mu}, \theta$ have unperturbed value zero, we 
avoid perturbations that merely represent coordinate transformation 
and have no physical significance.

To the first order the equations (1)–(4) and (7)–(9) are
\begin{align}
E_{ab}{}^{;b} &= \tfrac{1}{3} h_{a}{}^{b} \mu_{,b}, \tag{13}\\
H_{ab}{}^{;b} &= (\mu + h)\, \omega_{a}, \tag{14}\\
\dot{E}_{ab} + E_{c(a}\theta + h^{c}{}_{(a}\eta_{b)de}u^{d}H^{e}{}_{c}{}^{;c} 
  &= -\tfrac{1}{2}(\mu + h)\,\sigma_{ab}, \tag{15}\\
\dot{H}_{ab} + H_{c(a}\theta - h^{c}{}_{(a}\eta_{b)de}u^{d}E^{e}{}_{c}{}^{;c} &= 0, \tag{16}\\
\dot{\theta} &= -\tfrac{1}{3}\theta^{2} + \dot{u}^{a}{}_{;a} - \tfrac{1}{2}(\mu + 3h), \tag{17}\\
\dot{\omega}_{ab} &= -\tfrac{2}{3}\omega_{ab}\theta + \dot{u}_{[a;p]} h^{p}{}_{b} h^{q}{}_{a}, \tag{18}\\
\dot{\sigma}_{ab} &= E_{ab} - \tfrac{2}{3}\sigma_{ab}\theta - \tfrac{1}{3}h_{ab}\dot{u}^{c}{}_{;c}
   + \dot{u}_{(p;q)} h^{p}{}_{a} h^{q}{}_{b}. \tag{19}
\end{align}
From these we see that perturbations of rotation or of $E_{ab}$ or $H_{ab}$ do 
not produce perturbations of the expansion or the density. Nor do 
perturbations of $E_{ab}$ and $H_{ab}$ produce rotational perturbations.  

\sectiontoc{The Undisturbed Metric}

Since in the unperturbed state the rotation and acceleration 
are zero, $u_{a}$ must be hypersurface orthogonal.

\[
u_{a} = \tau_{,a},
\]

where $\tau$ measures the proper time along the world lines. As the 
surfaces $\tau = \text{constant}$ are homogeneous and isotropic they must be 
3-surfaces of constant curvature. Therefore the metric can be 
written,

\[
ds^{2} = - d\tau^{2} + \Omega^{2} d\gamma^{2},
\]

where  

\[
\Omega = \Omega(\tau),
\]

\[
d\gamma^{2} \quad \text{is the line element of a space of zero or unit positive or negative curvature.}
\]

We define $t$ by  

\[
\frac{dt}{d\tau} = \frac{1}{\Omega},
\]

then  

\[
ds^{2} = \Omega^{2} \left( -dt^{2} + d\gamma^{2} \right).
\]

In this metric,  

\[
u_{a} = (-\Omega, 0, 0, 0),
\]

\[
\therefore \quad \theta = \frac{3 \dot{\Omega}}{\Omega} = \frac{3 \Omega'}{\Omega^{2}}
\]

\[
\text{(prime denotes differentiation with respect to $t$).}
\]
Then, by (5), (7)
\begin{align}
\dot{\mu} &= -(\mu + h)\,\frac{3\dot{\Omega}}{\Omega}, \tag{20}\\[6pt]
3\frac{\ddot{\Omega}}{\Omega} &= -\tfrac{1}{2}(\mu + 3h). \tag{21}
\end{align}

If we know the relation between $\mu$ and $h$, we may determine $\Omega$.  
We will consider the two extreme cases, $h=0$ (dust) and $h=\tfrac{1}{3}\mu$ (radiation).  
Any physical situation should lie between these.

\subsubsection*{For $h = 0$}

By (20), \quad $\mu = \dfrac{M}{\Omega^{3}}$, \quad $M = \text{const}.$

\[
\therefore \quad \frac{3}{M}\,\Omega \ddot{\Omega} - \tfrac{1}{2}\Omega^{2} = 0,
\]

\[
\therefore \quad \frac{3}{M}\,\dot{\Omega}^{2} - \frac{1}{\Omega} = E, 
\qquad E = \text{const}.
\]

\paragraph{(a) For $E>0$,}
\[
\Omega = \tfrac{1}{2E}\Big(\cosh \sqrt{\tfrac{E M}{3}}\,t - 1\Big), 
\qquad 
\tau = \tfrac{1}{2E}\Big(\sqrt{\tfrac{3}{E M}}\,\sinh \sqrt{\tfrac{E M}{3}}\,t + t \Big);
\]

\paragraph{(b) For $E=0$,}
\[
\Omega = \frac{M}{12}\,t^{2}, 
\qquad 
\tau = \frac{M}{36}\,t^{3};
\]

\paragraph{(c) For $E<0$,}
\[
\Omega = -\tfrac{1}{2E}\Big(1 - \cos \sqrt{-\tfrac{E M}{3}}\,t \Big), 
\qquad 
\tau = -\tfrac{1}{2E}\Big(t - \sqrt{\tfrac{3}{-E M}} \sin \sqrt{-\tfrac{E M}{3}}\,t \Big).
\]

\medskip
$E$ represents the energy (kinetic + potential) per unit mass.  
If it is non–negative the universe will expand indefinitely, otherwise it will eventually contract again.
By the Gauss–Codazzi equations $^{*}R$, the curvature of the 
hypersurface $\tau = \text{const.}$ is  

\[
^{*}R = 2\left(-\tfrac{1}{3}\theta^{2} + \mu\right)
      = -\frac{2 E M}{\Omega^{2}}.
\]

If $E > 0$, \quad $^{*}R = -\dfrac{6}{\Omega^{2}}, \qquad M = \dfrac{3}{E}$;

\[
E = 0, \quad {}^{*}R = 0;
\]

If \(E < 0\), \quad
\({}^{*}R = \dfrac{6}{\Omega^{2}}, \qquad M = -\dfrac{3}{E}\).


\subsubsection*{For $h = \tfrac{1}{3}\mu$}

\[
\dot{\mu} = -4\frac{\dot{\Omega}}{\Omega},
\]

\[
3\frac{\ddot{\Omega}}{\Omega} = -\mu, 
\qquad \mu = \frac{M}{\Omega^{4}},
\]

\[
\therefore \quad \frac{3}{M}\dot{\Omega}^{2} - \frac{1}{\Omega^{2}} = E.
\]

\paragraph{(a) For $E > 0$:}
\[
\Omega = \frac{1}{E}\sinh t, 
\qquad \tau = \frac{1}{E}(\cosh t - 1), 
\qquad {}^{*}R = -\frac{6}{\Omega^{2}};
\]

\paragraph{(b) For $E = 0$:}
\[
\Omega = t, 
\qquad \tau = \tfrac{1}{2}t^{2}, 
\qquad {}^{*}R = 0;
\]

\paragraph{(c) For $E < 0$:}
\[
\Omega = \frac{1}{E}\sin t, 
\qquad \tau = \frac{1}{E}(\cos t - 1), 
\qquad {}^{*}R = \frac{6}{\Omega^{2}}.
\]

\sectiontoc{Rotational Perturbations}

By (6)
\[
\dot{u}_{[c;d]} h^{c}{}_{a} h^{d}{}_{b} 
   = -\frac{\omega_{ab}\,\dot{h}}{\mu + h}.
\]
\[
\dot{\omega}_{ab} = -\omega_{ab}\left(\tfrac{2}{3}\theta + \tfrac{\dot{h}}{\mu + h}\right).
\]

For $h = 0$,
\[
\omega = \frac{\omega_{0}}{\Omega^{2}}.
\]

For $h = \tfrac{\mu}{3}$,
\[
\dot{\omega} = -\omega\left(\tfrac{2}{3}\theta + \tfrac{1}{4}\frac{\dot{\mu}}{\mu}\right)
= -\tfrac{5}{3}\omega\theta,
\]
\[
\therefore \quad \omega = \frac{\omega_{0}}{\Omega^{5}}.
\]

Thus rotation dies away as the universe expands.  
This is in fact a statement of the conservation of angular momentum in an expanding universe.

\sectiontoc{Perturbations of Density}

For $h = 0$ we have the equations,
\[
\dot{\mu} = -\mu\theta,
\]
\[
\dot{\theta} = -\tfrac{1}{3}\theta^{2} - \tfrac{1}{2}\mu.
\]

These involve no spatial derivatives. Thus the behaviour of one 
region is unaffected by the behaviour of another. Perturbations 
will consist in some regions having slightly higher or lower values 
of $\theta$ than the average. If the universe as whole has a value of $E$ 
greater than zero, a small perturbation will still have $E$ greater 
than zero and will continue to expand. It will not contract to form a galaxy. If the universe has a value of $E$ less than zero, a 
small perturbation can contract. However it will only begin 
contracting at a time $\delta\tau$ earlier than the whole universe begins 
contracting, where  

\[
\frac{\delta\tau}{\tau_{c}} = \frac{\delta E}{E_{0}},
\]

$\tau_{c}$ is the time at which the whole universe begins contracting.  
There is only any real instability when $E=0$. This case is of 
measure zero relative to all the possible values $E$ can have.  
However this cannot really be used as an argument to dismiss it 
as there might be some reason why the universe should have $E=0$.

For a region with energy $- \delta E$, in a universe with $E=0$,

\[
\Omega = \frac{1}{4\delta E}\left(t^{2} - \frac{t^{4}}{12} + \cdots \right),
\]

\[
\tau = \frac{1}{12\delta E}\left(t^{3} - \frac{t^{5}}{20} + \cdots \right),
\]

\[
\mu = \frac{3}{6E\,\Omega^{2}} = \frac{4}{3}\tau^{-2}
     \left(1 + \frac{(\delta E)^{2/3}}{2^{1/3}}\,\tau^{2/3} + \cdots \right).
\]

For $E=0$, \quad $\mu = \tfrac{4}{3}\tau^{-2}$.

Thus the perturbation grows only as $\tau^{2/3}$.  
This is not fast enough to produce galaxies from statistical fluctuations 
even if these could occur. However, since an evolutionary universe has 
a particle horizon (Rindler\,(8), Penrose\,(9)) different parts do not 
communicate in the early stages. This makes it even more difficult for 
statistical fluctuations to occur over a region until light had time 
to cross the region.

\subsubsection*{For $h = \tfrac{\mu}{3}$}

\[
\dot{\mu} = -\tfrac{4}{3}\mu\theta,
\]
\[
\dot{\theta} = -\tfrac{1}{3}\theta^{2} - \mu + \dot{u}^{a}{}_{;a},
\]
\[
\dot{u}_{a} = -\frac{h^{b}{}_{a}\,\mu_{,b}}{4\mu}.
\]

As before, a perturbation cannot contract unless it has a negative 
value of $E$. The action of the pressure forces makes it still more 
difficult for it to contract. Eliminating $\theta$,

\[
\mu \ddot{\mu} - \tfrac{5}{4}\dot{\mu}^{2} - \tfrac{5}{3}\mu^{3} + \tfrac{4}{3}\mu^{2}\dot{u}^{a}{}_{;a} = 0,
\]

\[
\ddot{u}^{a} = \dot{u}^{a}{}_{;b}\,h^{b}{}_{a}{}^{b} + \dot{u}^{a}\dot{u}^{a}
= -\tfrac{1}{4}\,h^{ac}\frac{(h^{b}{}_{a}\mu_{,b})_{;c}}{\mu}
\quad \text{to our approximation.}
\]

\[
h^{ac}\nabla_{c} h^{b}{}_{a}\nabla_{b}
\]
is the Laplacian in the hypersurface $\tau = \text{constant}$.  

We represent the perturbation as a sum of eigenfunctions $S^{(n)}$ of 
this operator, where  

\[
S^{(n)}{}_{;a}u^{a} = 0,
\]

\[
h^{ac}(h^{b}{}_{a}S^{(n)}{}_{,b})_{;c} 
= -\frac{n^{2}}{\Omega^{2}}\,S^{(n)}.
\]

These eigenfunctions will be hyperspherical and pseudohyperspherical 
harmonics in cases (c) and (a) respectively and plane waves in case (b).  
In case (c) $n$ will take only discrete values but in (a) and (b) it will take all positive values.

\[
\mu = \mu_{0}\left(1 + \sum_{n} B^{(n)} S^{(n)}\right),
\]

where $\mu_{0}$ is the undisturbed density.

\[
\ddot{B}^{(n)}\mu_{0} - \tfrac{1}{2}\dot{B}^{(n)}\mu_{0} - B^{(n)}\left(\tfrac{4}{3}\mu_{0}^{2} - \tfrac{n^{2}}{3\Omega^{2}}\mu_{0}\right) = 0.
\]

As long as $\mu_{0} > \tfrac{n^{2}}{4\Omega^{2}}$, $B^{(n)}$ will grow.  

For $\mu_{0} \gg \tfrac{n^{2}}{4\Omega^{2}}$,

\[
B^{(n)} \simeq C\,\tau + D\,\tau^{-1}.
\]

These perturbations grow for as long as light has not had time to 
travel a significant distance compared to the scale of the perturbation 
($\sim \tfrac{\Omega}{n}$). Until that time pressure forces cannot act to even out 
perturbations.

When $\tfrac{n^{2}}{\Omega^{2}} \gg \mu_{0}$,

\[
\ddot{B}^{(n)} + \dot{B}^{(n)}\frac{\dot{\Omega}}{\Omega} + \tfrac{n^{2}}{3}B^{(n)} = 0,
\]

\[
\therefore \quad B^{(n)} \simeq C\,\Omega^{-\tfrac{1}{2}} e^{\pm i \tfrac{n}{\sqrt{3}} t}.
\]

We obtain sound waves whose amplitude decreases with time.  
These results confirm those obtained by Lifshitz and Khalatnikov\,(3).

From the foregoing we see that galaxies cannot form as the result 
of the growth of small perturbations. We may expect that other non–
gravitational forces will have an effect smaller than pressure equal 
\sectiontoc{The Steady-state Universe}

To obtain the steady-state universe we must add extra terms to 
the energy–momentum tensor. Hoyle and Narlikar\,(10) use,

\[
T_{ab} = \mu\, u_{a} u_{b} + h_{ab} - C_{a} C_{b} 
+ \tfrac{1}{2} g_{ab} C_{c} C^{d}. \tag{20}
\]

where,  

\[
C_{a} = C_{;a}, \qquad 
C_{;a}{}^{a} = -\dot{j}_{a}{}^{a}, \qquad 
j_{a} = -(\mu + h) u_{a}.
\]

Since $T_{ab}{}^{;b} = 0$,

\[
\dot{\mu} + (\mu + h)\theta + u^{a} C_{\mu} C_{b}{}^{;b} = 0, \tag{21}
\]

\[
(\mu + h)\dot{u}_{a} + h^{b}{}_{a} h^{c}{}_{b} 
- h^{b}{}_{a} C_{b} C_{d}{}^{;d} = 0. \tag{22}
\]

There is a difficulty here, if we require that the ``$C$'' field


should not produce acceleration or, in other words, that the matter 
created should have the same velocity as the matter already in 
existence. We must then have

\[
h^{b}{}_{a} C_{b} = 0. \tag{23}
\]

However since $C$ is a scalar, this implies that the rotation of the 
medium is zero. On the other hand if (23) does not hold, the equations 
are indeterminate (c.f. Raychaudhuri and Bannerjee\,(11)). In order to 
have a determinate set of equations we will adopt (23) but drop the 
requirement that $C_{a}$ is the gradient of a scalar. The condition (23) 
is not very satisfactory but it is difficult to think of one more 
satisfactory. Hoyle and Narlikar\,(12) seek to avoid this difficulty 
by taking a particle rather than a fluid picture. However this has a 
serious drawback since it leads to infinite fields (Hawking\,(13)).

From (17),

\[
C_{a} = -u_{a}\left[1 - \frac{\dot{\mu}}{\mu + \dot{\mu} + (\mu + h)\theta}\right].
\]

\[
\therefore \quad C^{a}{}_{;a} = -(\dot{\mu}+h) - (\mu+h)\theta,
\]

\[
= -\theta\left(1 - \frac{\dot{\mu}}{\mu + \dot{\mu} + (\mu + h)\theta}\right)
+ \frac{\dot{\mu}}{\mu + \dot{\mu} + (\mu + h)\theta}.
\]


For $\mu \gg h$
\[
(\dot{\mu} + \dot{h}) = \theta \big(1 - (\mu + h)\big),
\]
\[
\therefore \quad (\mu + h) \to 1.
\]

Thus, small perturbations of density die away. Moreover equation (18) 
still holds, and therefore rotational perturbations also die away.  
Equation (19) now becomes

\[
\dot{\theta} = -\tfrac{1}{3}\theta^{2} - \tfrac{1}{2}(\mu + 3h) + 1,
\]
\[
\therefore \quad \theta \to \sqrt{3\left(\tfrac{1}{2} - h\right)}.
\]

These results confirm those obtained by Hoyle and Narlikar\,(14). We 
see therefore that galaxies cannot be formed in the steady-state 
universe by the growth of small perturbations. However this does not 
exclude the possibility that there might be a self-perpetuating 
system of finite perturbations which could produce galaxies.  
(Sciama\,(15), Roxburgh and Saffman\,(16)).

\sectiontoc{Gravitational Waves}

We now consider perturbations of the Weyl tensor that do not 
arise from rotational or density perturbations, that is,
\[
E_{ab}{}^{;b} = H_{ab}{}^{;b} = 0
\]

Multiplying (15) by $u^{c}\nabla_{c}$, and (16) by 
$h^{a}{}_{(c}\eta_{b)}{}^{rs}u_{r}\nabla_{s}$, 
we obtain, after a lot of reduction,
\[
\ddot{E}_{ab} - \left( E_{cd;e}\, h^{c}{}_{f} h^{d}{}_{g} h^{e}{}_{(a}\, h^{fg}{}_{b)} \right)
+ h^{kl} h_{a}{}^{t} h_{b}{}^{s}\, \tfrac{7}{3}\,\dot{E}_{ab}\,\theta
\]
\[
+ E_{ab}\left(\dot{\theta} + \tfrac{4}{3}\theta^{2} + \tfrac{1}{3}(\mu + 3h)\right) 
+ \sigma_{ab}\left(\tfrac{1}{3}\theta(\mu+h) + \tfrac{1}{2}(\dot{\mu}+\dot{h})\right) = 0. \tag{24}
\]

In empty space with a non-expanding congruence $u^{a}$ this reduces to 
the usual form of the linearised theory,
\[
\square^{*} E_{ab} = 0.
\]

The second term in (24) is the Laplacian in the hypersurface 
$\tau = \text{constant}$, acting on $E_{ab}$. We will write $E_{ab}$ as a sum of 
eigenfunctions of this operator,
\[
E_{ab} = \sum A^{(n)} V_{ab}^{(n)},
\]
where
\[
\dot{V}_{ab}^{(n)} = 0, \qquad
\left( V_{cd;e}\, h^{c}{}_{f} h^{d}{}_{g} h^{e}{}_{(a} h^{fg}{}_{b)} \right)
= -\frac{n^{2}}{\Omega^{2}}\, V_{ab}^{(n)},
\]
\[
V_{ab}{}^{;b} = 0, \qquad V_{a}{}^{a} = 0.
\]

Then
\[
\dot{E}_{ab} = \sum \frac{A^{(n)}}{\Omega}\, V_{ab}^{(n)},
\qquad
\ddot{E}_{ab} = \sum \left[\frac{\dot{A}^{(n)}}{\Omega^{2}} 
- \frac{A^{(n)}}{\Omega^{2}}\right] V_{ab}^{(n)}.
\]

Similarly,
\[
\sigma_{ab} = \sum D^{(n)} V_{ab}^{(n)}.
\]

Then by (19)
\[
D'^{(n)} = \Omega A^{(n)} - 2D^{(n)}\frac{\Omega'}{\Omega}.
\]

Substituting in (24),
\[
A''^{(n)} + \frac{6\Omega'}{\Omega} A'^{(n)} + A^{(n)} 
\left[ n^{2} + 3\frac{\Omega''}{\Omega} + 6\frac{\Omega'^{2}}{\Omega^{2}} 
+ \tfrac{1}{3}(\mu + 3h)\Omega^{2}\right]
+ D^{(n)}\left[\Omega(\mu+h) + \tfrac{1}{2}\Omega^{2}(\dot{\mu}+\dot{h})\right] = 0.
\]

We may differentiate again and substitute for $D'$.  
For $n \gg 1$ and $\Omega \gg \tfrac{n^{2}}{h^{2}}$,
\[
A^{(n)} \simeq \frac{1}{\Omega^{3/2}} e^{i n t}.
\]

So the gravitational field $E_{ab}$ decreases as $\Omega^{-1}$ and the 
``energy'' $\tfrac{1}{2}(E_{ab}E^{ab} + H_{ab}H^{ab})$ as $\Omega^{-6}$.  
We might expect this as the Bianchi identities may be written, to the 
linear approximation,
\[
\Omega\, g^{ed} \frac{\partial}{\partial x^{e}}
\left(\Omega^{2} C_{abcd}\right) = J_{abc}.
\]

Therefore if the interaction with the matter could be neglected 
$C_{abcd}$ would be proportional to $\Omega$ and $E_{ab}, H_{ab}$ to $\Omega^{-1}$.

In the steady-state universe when $\mu$ and $\theta$ have reached their 
equilibrium values,
\[
R_{ab} = \left(\tfrac{1}{2}\mu + h\right) g_{ab},
\]

\[
J_{abc} = R_{c[a;b]} - \tfrac{1}{6} g_{c[a}R_{;b]} = 0.
\]
Thus the interaction of the ``$C$'' field with gravitational radiation is 
equal and opposite to that of the matter. There is then no net 
interaction, and $E_{ab}$ and $H_{ab}$ decrease as $\Omega^{-1}$.

The ``energy'' $\tfrac{1}{2}(E_{ab}E^{ab} + H_{ab}H^{ab})$ depends on second derivatives 
of the metric. It is therefore proportional to the frequency squared 
times the energy as measured by the energy momentum pseudo-tensor, in 
a local co-moving Cartesian coordinate system, which depends only on 
first derivatives. Since the frequency will be inversely proportional 
to $\Omega$, the energy measured by the pseudo-tensor will be proportional 
to $\Omega^{-4}$ as for other rest mass zero fields.

\sectiontoc{Absorption of Gravitational Waves}

As we have seen, gravitational waves are not absorbed by a 
perfect fluid. Suppose however there is a small amount of viscosity.  
We may represent this by the addition of a term $\lambda\sigma_{ab}$ to the 
energy–momentum tensor, where $\lambda$ is the coefficient of viscosity 
(Ehlers\,(17)).

Since
\[
T_{ab}{}^{;b} = 0
\]
we have
\[
\dot{\mu} + (\mu+h)\theta - 2\lambda\sigma^{2} = 0, \tag{25}
\]
\[
(\mu+h)\dot{u}_{a} + h^{b}{}_{a}\mu_{,b} + \lambda\sigma_{cb}{}^{;b}h^{c}{}_{a} = 0. \tag{26}
\]
Equations (15) (16) become
\[
\dot{E}_{ab} + E_{c(a}\theta + h^{c}{}_{(a}\eta_{b)de}u^{c} H_{f}{}^{d;e}
= -\tfrac{1}{2}(\mu + h)\sigma_{ab}
- \tfrac{1}{2}\lambda\left(E_{ab} - \tfrac{1}{3}\sigma_{ab}\theta\right), \tag{27}
\]

\[
\dot{H}_{ab} + H_{c(a}\theta - h^{c}{}_{(a}\eta_{b)de}u^{c} E_{f}{}^{d;e}
= -\tfrac{1}{2}\lambda H_{ab}. \tag{28}
\]

The extra terms on the right of equations (27), (28) are similar to 
conduction terms in Maxwell’s equations and will cause the wave to 
decrease by a factor $e^{-\tfrac{1}{2}\lambda t}$. Neglecting expansion for the moment, 
suppose we have a wave of the form,
\[
E_{ab} = E^{0}_{ab} e^{i\nu\tau}.
\]

This will be absorbed in a characteristic time $2/\lambda$ independent of 
frequency. By (25) the rate of gain of rest mass energy of the 
matter will be $2\lambda\sigma^{2}$ which by (19) will be $2\lambda E^{2}\nu^{-2}$. 
Thus the available energy in the wave is $4 E^{2}\nu^{2}$. This confirms that the 
density of available energy of gravitational radiation will decrease 
as $\Omega^{-4}$ in an expanding universe. From this we see that 
gravitational radiation behaves in much the same way as other 
radiation fields. In the early stages of an evolutionary universe 
when the temperature was very high we might expect an equilibrium to 
be set up between black-body electromagnetic radiation and black-body 
gravitational radiation. Since they both have two polarisations their energy densities should be equal. As the universe expanded they would 
both cool adiabatically at the same rate. As we know the temperature 
of black-body extragalactic electromagnetic radiation is less than 
$5^\circ$K, the temperature of the black-body gravitational radiation must 
be also less than this which would be absolutely undetectable. Now 
the energy of gravitational radiation does not contribute to the 
ordinary energy–momentum tensor $T_{ab}$. Nevertheless it will have an 
active gravitational effect. By the expansion equation,
\[
\dot{\theta} = -\tfrac{1}{3}\theta^{2} - 2\sigma^{2} - \tfrac{1}{2}(\mu + 3h).
\]

For incoherent gravitational radiation at frequency $\nu$,
\[
\sigma^{2} = \alpha E^{2}\nu^{-2}.
\]

But the energy density of the radiation is
\[
4\,E^{2}\nu^{-2}.
\]

\[
\therefore \quad \dot{\theta} = -\tfrac{1}{3}\theta^{2} - \tfrac{1}{2}\mu_{G} - \tfrac{1}{2}(\mu + 3h),
\]
where $\mu_{G}$ is the gravitational ``energy'' density. Thus gravitational 
radiation has an active attractive gravitational effect. It is 
interesting that this seems to be just half that of electromagnetic 
radiation.

It has been suggested by Hogarth\,(18) and Hoyle and Narlikar\,(10), 
that there may be a connection between the absorption of radiation 
and the Arrow of Time. Thus in universes like the steady-state, in 
which all electromagnetic radiation emitted is eventually absorbed by 
other matter, the Absorber theory would predict retarded solutions of the Maxwell equations while in evolutionary universes in which 
electromagnetic radiation is not completely absorbed it would predict 
advanced solutions. Similarly, if one accepted this theory, one would 
expect retarded solutions of the Einstein equations if and only if all 
gravitational radiation emitted is eventually absorbed by other matter. 
Clearly this is so for the steady–state universe since $\lambda$ will be 
constant. In evolutionary universes $\lambda$ will be a function of time.  
We will obtain complete absorption if $\int \lambda d\tau$ diverges. Now for a gas, 
$\lambda \propto T^{2}$ where $T$ is the temperature. For a monatomic gas, $T \propto \Omega^{-2}$, 
therefore the integral will diverge (just). However the expression 
used for viscosity assumed that the mean free path of the atoms was 
small compared to the scale of the disturbance. Since the mean free 
path $\propto \mu^{-1}\kappa \Omega^{-3}$ and the wavelength $\propto \Omega^{-1}$, the mean free path will 
eventually be greater than the wavelength and so the effective viscosity 
will decrease more rapidly than $\Omega^{-1}$. Thus there will not be complete 
absorption and the theory would not predict retarded solutions.  

However this is slightly academic since gravitational radiation has not 
yet been detected, let alone investigated to see whether it corresponds 
to a retarded or advanced solution.

\sectiontoc{References}

\begin{enumerate}
\item Bonnor W.B.; M.N., \textbf{117}, 104 (1957).
\item Lifshitz E.M.; J.Phys. U.S.S.R., \textbf{10}, 116 (1946).
\item Lifshitz E.M., Khalatnikov I.N.; Adv. in Phys., \textbf{12}, 185 (1963).
\item Irvine W.; Annals of Physics, \textbf{32}, 322 (1965).
\item Jordan P., Ehlers J., Kundt W.; Ahb.Akad.Wiss. Mainz., No.7 (1960).
\item Kundt W., Trümper M.; Ahb.Akad.Wiss. Mainz., No.12 (1962).
\item Trümper M.; Contributions to Actual Problems in Gen.Rel.
\item Rindler W.; M.N., \textbf{116}, 662 (1956).
\item Penrose R.; Article in `Relativity Topology and Groups’, Gordon and Breach (1964).
\item Hoyle F., Narlikar J.V.; Proc.Roy.Soc., A, \textbf{277}, 1 (1964).
\item Raychaudhuri A., Banerjee S.; Z.Astrophys., \textbf{58}, 187 (1964).
\item Hoyle F., Narlikar J.V.; Proc.Roy.Soc., A, \textbf{282}, 184 (1964).
\item Hawking S.W.; Proc.Roy.Soc., A, \textbf{286}, 313 (1965).
\item Hoyle F., Narlikar J.V.; Proc.Roy.Soc., A, \textbf{273}, 1 (1963).
\item Sciama D.W.; M.N., \textbf{115}, 3 (1955).
\item Roxburgh I.W., Sattman P.G.; M.N., \textbf{129}, 181 (1965).
\item Ehlers J.; Ahb.Akad.Wiss. Mainz. Noll. (1961).
\item Hogarth J.E.; Proc.Roy.Soc., A, \textbf{267}, 365 (1962).
\end{enumerate}




